\documentclass[a4paper,11pt]{article}

\usepackage{fontspec}
\usepackage{libertinus-otf}
\usepackage[french]{babel}

\usepackage{amsmath,amssymb,amsfonts}
\usepackage{array}
\usepackage{booktabs}
\usepackage{longtable}
\usepackage{geometry}
\usepackage{multirow}
\usepackage[table]{xcolor}
\usepackage{colortbl}
\usepackage{fancyhdr}
\usepackage{lastpage}
\usepackage{enumitem}
\usepackage{eurosym}
\usepackage{tikz}
\usepackage[most]{tcolorbox}

\geometry{
    a4paper,
    left=1.1cm,
    right=1.1cm,
    top=1.6cm,
    bottom=1.5cm,
    headheight=20pt,
    headsep=10pt,
    footskip=20pt
}

\setlength{\parindent}{0pt}
\setlength{\tabcolsep}{4.5pt}
\renewcommand{\arraystretch}{1.18}
\setlength{\aboverulesep}{1.5pt}
\setlength{\belowrulesep}{1.5pt}

\definecolor{bleuprimaire}{RGB}{20, 52, 98}
\definecolor{bleusecondaire}{RGB}{31, 95, 168}
\definecolor{bleufond}{RGB}{238, 244, 251}
\definecolor{grisaccent}{RGB}{100, 116, 139}
\definecolor{vertvrai}{RGB}{16, 120, 60}
\definecolor{rougefaux}{RGB}{195, 20, 20}
\definecolor{filetgris}{RGB}{220, 226, 235}

\tikzset{
    tabcadre/.style={draw=bleuprimaire, line width=0.8pt},
    tabligne/.style={draw=bleuprimaire!45, line width=0.6pt},
    tabrep/.style={
        draw=rougefaux!90,
        fill=rougefaux!8,
        text=rougefaux,
        font=\bfseries\footnotesize,
        rounded corners=2.5pt,
        inner sep=1.8pt,
        line width=0.6pt
    },
    tabfleche/.style={
        ->,
        >=stealth,
        thick,
        draw=bleusecondaire,
        shorten <= 3pt,
        shorten >= 3pt
    }
}

\newcommand{\pts}[1]{%
    \tcbox[on line, size=small, colback=bleufond, colframe=bleusecondaire!50,
           arc=2.5mm, boxrule=0.6pt, top=1.2pt, bottom=1.2pt, left=4pt, right=4pt]{%
        \small\bfseries\textcolor{bleuprimaire}{#1}%
    }%
}

\newcommand{\repbox}[1]{%
    \tcbox[on line, size=small, colback=rougefaux!6, colframe=rougefaux!90,
           arc=1.2mm, boxrule=0.7pt, top=1pt, bottom=1pt, left=3.5pt, right=3.5pt]{%
        \bfseries\textcolor{rougefaux}{#1}%
    }%
}

\newcommand{\vrai}{\textcolor{vertvrai}{\textbf{[VRAI]}}}
\newcommand{\faux}{\textcolor{rougefaux}{\textbf{[FAUX]}}}
\newcommand{\subpt}[1]{\textcolor{bleusecondaire}{\bfseries [#1]}}
\newcommand{\tol}[1]{\textcolor{grisaccent}{\footnotesize\textit{(#1)}}}

\pagestyle{fancy}
\fancyhf{}
\lhead{\small\bfseries\textcolor{bleuprimaire}{BAC BLANC} \textcolor{grisaccent}{-- Terminale Spécialité Mathématiques}}
\rhead{\small\bfseries\textcolor{bleusecondaire}{Corrigé et Barème Officiel}}
\cfoot{\small\textcolor{grisaccent}{Page \thepage\ sur \pageref{LastPage}}}

\renewcommand{\headrulewidth}{0.5pt}
\renewcommand{\headrule}{\hbox to\headwidth{\color{bleusecondaire!60}\leaders\hrule height \headrulewidth\hfill}}
\renewcommand{\footrulewidth}{0.3pt}
\renewcommand{\footrule}{\hbox to\headwidth{\color{filetgris}\leaders\hrule height \footrulewidth\hfill}}

\begin{document}

\begin{tcolorbox}[
    enhanced,
    colback=bleufond,
    colframe=bleuprimaire,
    arc=1.6mm,
    boxrule=0.8pt,
    top=4pt, bottom=4pt, left=8pt, right=8pt
]
    \textbf{\Large\textcolor{bleuprimaire}{Grille d'évaluation chiffrée et Corrigé détaillé}} \hfill \textbf{\Large\textcolor{bleuprimaire}{Total : 20 points}} \\[0.15em]
    \textbf{\textcolor{bleusecondaire}{Sujet Unique}} \hspace{0.2cm}\textcolor{grisaccent}{\textbullet}\hspace{0.2cm} \textcolor{grisaccent}{Bienaymé-Tchebychev & Équations Différentielles} \hfill \small\textcolor{grisaccent}{Durée : 4 heures}
\end{tcolorbox}

\vspace{0.10cm}

\begin{longtable}{
    >{\centering\arraybackslash\bfseries}p{1.1cm}
    >{\raggedright\arraybackslash\bfseries}p{4.1cm}
    !{\color{filetgris}\vrule width 0.7pt}
    >{\raggedright\arraybackslash}p{11.1cm}
    >{\centering\arraybackslash}p{1.2cm}
}

\toprule[1.2pt]
\rowcolor{bleuprimaire}
\textcolor{white}{\bfseries Réf.} &
\textcolor{white}{\bfseries Compétence / Question} &
\textcolor{white}{\bfseries Éléments de correction & Barème détaillé} &
\textcolor{white}{\bfseries Barème} \\
\midrule[1pt]
\endfirsthead

\toprule[1.2pt]
\rowcolor{bleuprimaire}
\textcolor{white}{\bfseries Réf.} &
\textcolor{white}{\bfseries Compétence / Question} &
\textcolor{white}{\bfseries Éléments de correction & Barème détaillé} &
\textcolor{white}{\bfseries Barème} \\
\midrule[1pt]
\endhead

% ==================== EXERCICE 1 ====================
\rowcolor{bleusecondaire!15}
\multicolumn{4}{l}{
    \bfseries\textcolor{bleuprimaire}{\large Exercice 1 \hspace{0.2cm}\textbullet\hspace{0.2cm} Probabilités -- Inégalité de Bienaymé-Tchebychev} \hfill \bfseries 4 points
} \\* \midrule[0.6pt]

\textbf{Q.1} &
Espérance et variance de $X$ et $Y_n$ &
\textbullet\ \subpt{0,50} $E(X)=0,4n$, $V(X)=0,24n$.\newline
\textbullet\ \subpt{0,50} $E(Y_n)=0,4$ et $V(Y_n) = \frac{0,24}{n}$ : \repbox{Formules exactes}. &
\pts{1,00} \\
\cmidrule(lr){2-4}
\textbf{Q.2} &
Énoncé Bienaymé-Tchebychev &
\textbullet\ \subpt{0,50} Rappel rigoureux de l'inégalité générale.\newline
\textbullet\ \subpt{0,50} Application aux variables aléatoires d'espérance finie : \repbox{Validé}. &
\pts{1,00} \\
\cmidrule(lr){2-4}
\textbf{Q.3} &
Démonstration pour $Y_n$ &
\textbullet\ \subpt{0,50} Substitution de $\mu$ et $\sigma^2$.\newline
\textbullet\ \subpt{0,50} Obtention de $\frac{0,24}{n\varepsilon^2}$ : \repbox{Démonstration complète}. &
\pts{1,00} \\
\cmidrule(lr){2-4}
\textbf{Q.4} &
Détermination du seuil $n$ &
\textbullet\ \subpt{0,50} Mise en inéquation $\frac{0,24}{n(0,05)^2} \leq 0,05$.\newline
\textbullet\ \subpt{0,50} Conclusion $n \geq 1920$ : \repbox{1920 épreuves}. &
\pts{1,00} \\

\midrule[1pt]

% ==================== EXERCICE 2 ====================
\rowcolor{bleusecondaire!15}
\multicolumn{4}{l}{
    \bfseries\textcolor{bleuprimaire}{\large Exercice 2 \hspace{0.2cm}\textbullet\hspace{0.2cm} Analyse -- Équations différentielles et Intégration} \hfill \bfseries 5 points
} \\* \midrule[0.6pt]

\textbf{A.1} &
Équation homogène $(E_0)$ &
\textbullet\ \subpt{0,50} Forme générale $y' + ay = 0$.\newline
\textbullet\ \subpt{0,50} Solution $y(x) = C\mathrm{e}^{-2x}$ : \repbox{C\mathrm{e}^{-2x}}. &
\pts{1,00} \\
\cmidrule(lr){2-4}
\textbf{A.2} &
Solution particulière $f_p$ &
\textbullet\ \subpt{0,50} Dérivation de $f_p(x) = ax\mathrm{e}^{-2x}$.\newline
\textbullet\ \subpt{0,50} Identification $a = 4$ : \repbox{4x\mathrm{e}^{-2x}}. &
\pts{1,00} \\
\cmidrule(lr){2-4}
\textbf{A.3} &
Ensemble des solutions de $(E)$ &
\textbullet\ \subpt{0,75} Somme solution homogène et particulière : \repbox{(C+4x)\mathrm{e}^{-2x}}. &
\pts{0,75} \\
\cmidrule(lr){2-4}
\textbf{A.4} &
Condition initiale $f(0) = 3$ &
\textbullet\ \subpt{0,75} Détermination de $C = 3$ : \repbox{(3+4x)\mathrm{e}^{-2x}}. &
\pts{0,75} \\
\cmidrule(lr){2-4}
\textbf{A.5} &
Limites, variations et intégrale $I$ &
\textbullet\ \subpt{0,50} Limite en $+\infty$ nulle.\newline
\textbullet\ \subpt{0,50} Variations et tableau.\newline
\textbullet\ \subpt{0,50} Calcul exact de l'intégrale $I$ : \repbox{Valeur exacte}. &
\pts{1,50} \\

\midrule[1pt]

% ==================== EXERCICE 3 ====================
\rowcolor{bleusecondaire!15}
\multicolumn{4}{l}{
    \bfseries\textcolor{bleuprimaire}{\large Exercice 3 \hspace{0.2cm}\textbullet\hspace{0.2cm} Nombres complexes et Géométrie} \hfill \bfseries 6 points
} \\* \midrule[0.6pt]

\textbf{A.1} &
Équation du second degré &
\textbullet\ \subpt{0,50} Calcul du discriminant $\Delta = -16$.\newline
\textbullet\ \subpt{0,75} Solutions complexes conjuguées $2 \pm 2\mathrm{i}$ : \repbox{2\pm 2\mathrm{i}}. &
\pts{1,25} \\
\cmidrule(lr){2-4}
\textbf{A.2} &
Modules, arguments et triangle $OAB$ &
\textbullet\ \subpt{0,75} $|z_A| = 2\sqrt{2}$, $\arg(z_A) = \frac{\pi}{4}$.\newline
\textbullet\ \subpt{0,75} Égalité $OA = OB = 2\sqrt{2}$ : \repbox{Isocèle en O}. &
\pts{1,50} \\
\cmidrule(lr){2-4}
\textbf{B.1} &
Transformation $f$ et image $A'$ &
\textbullet\ \subpt{0,75} Nature : similitude directe (rotation + translation).\newline
\textbullet\ \subpt{0,50} Affixe de $A'$ : \repbox{-2+4\mathrm{i}}. &
\pts{1,25} \\
\cmidrule(lr){2-4}
\textbf{B.2} &
Suite complexe et convergence &
\textbullet\ \subpt{0,50} Calculs de $z_1$ et $z_2$.\n\newline
\textbullet\ \subpt{0,75} Suite géométrique de raison $\frac{1}{2}\mathrm{i}$.\newline
\textbullet\ \subpt{0,75} Expression générale et limite : \repbox{Convergence}. &
\pts{2,00} \\

\midrule[1pt]

% ==================== EXERCICE 4 ====================
\rowcolor{bleusecondaire!15}
\multicolumn{4}{l}{
    \bfseries\textcolor{bleuprimaire}{\large Exercice 4 \hspace{0.2cm}\textbullet\hspace{0.2cm} Analyse -- Étude de fonction et Intégration par parties} \hfill \bfseries 5 points
} \\* \midrule[0.6pt]

\textbf{A.1} &
Limite en $0^+$ de $g$ &
\textbullet\ \subpt{0,50} Croissance comparée $x\ln(x) \to 0$.\newline
\textbullet\ \subpt{0,50} Limite globale : \repbox{0}. &
\pts{1,00} \\
\cmidrule(lr){2-4}
\textbf{A.2} &
Variations de $g$ &
\textbullet\ \subpt{0,75} Dérivée $g'(x) = \ln(x)$.\newline
\textbullet\ \subpt{0,75} Tableau de variations complet : \repbox{Minimum en 1}. &
\pts{1,50} \\
\cmidrule(lr){2-4}
\textbf{A.3} &
Intégrale $J$ par IPP &
\textbullet\ \subpt{1,00} Intégration par parties rigoureuse : \repbox{J = 1}. &
\pts{1,00} \\
\cmidrule(lr){2-4}
\textbf{B.1} &
Primitive et calcul d'aire &
\textbullet\ \subpt{0,75} Dérivation de $h$ vérifiant $h'(x) = g(x)$.\newline
\textbullet\ \subpt{0,75} Calcul de l'aire exacte : \repbox{Unités d'aire}. &
\pts{1,50} \\

\midrule[1.2pt]
\rowcolor{bleuprimaire}
\multicolumn{3}{r}{\textcolor{white}{\bfseries\large TOTAL GÉNÉRAL}} &
\textcolor{white}{\bfseries\large 20 pts} \\
\bottomrule[1.2pt]

\longtable}

\end{document}
